\section*{Sets and Indices}

No nontrivial index sets are used in the implemented model, since there are only two vehicle types represented by separate scalar decision variables.

\section*{Parameters}

\[
S = 400
\qquad\text{(code data: \texttt{num\_students})}
\]
Total number of students to be transported.

\[
B^{\max} = 10
\qquad\text{(code data: \texttt{num\_buses})}
\]
Maximum number of buses available.

\[
Q^B = 50
\qquad\text{(code data: \texttt{bus\_capacity})}
\]
Capacity of one bus.

\[
M^{\max} = 8
\qquad\text{(code data: \texttt{num\_minibuses})}
\]
Maximum number of minibuses available.

\[
Q^M = 40
\qquad\text{(code data: \texttt{minibus\_capacity})}
\]
Capacity of one minibus.

\[
D = 9
\qquad\text{(code data: \texttt{num\_drivers})}
\]
Total number of available drivers.

\[
c^B = 800
\qquad\text{(code data: \texttt{bus\_rental\_cost})}
\]
Rental cost of one bus.

\[
c^M = 600
\qquad\text{(code data: \texttt{minibus\_rental\_cost})}
\]
Rental cost of one minibus.

\section*{Decision Variables}

\[
x \in \mathbb{Z}
\qquad\text{(code: \texttt{x})}
\]
Number of buses rented/used.

\[
y \in \mathbb{Z}
\qquad\text{(code: \texttt{y})}
\]
Number of minibuses rented/used.

The code also imposes bounds
\[
0 \le x \le B^{\max}, \qquad 0 \le y \le M^{\max}.
\]

\section*{Objective}

Minimize total vehicle rental cost:
\[
\min \; c^B x + c^M y .
\]

For the given instance,
\[
\min \; 800x + 600y.
\]

\section*{Constraints}

\paragraph{Student transportation capacity.}
The total seating capacity of the selected vehicles must cover all students:
\[
Q^B x + Q^M y \ge S.
\]
For the given instance,
\[
50x + 40y \ge 400.
\]

\paragraph{Driver limit.}
Each vehicle requires one driver, and the total number of drivers is limited:
\[
x + y \le D.
\]
For the given instance,
\[
x + y \le 9.
\]

\paragraph{Vehicle availability bounds.}
The number of selected vehicles cannot exceed availability:
\[
0 \le x \le B^{\max},
\]
\[
0 \le y \le M^{\max}.
\]

\paragraph{Integrality.}
Vehicle counts must be integer:
\[
x, y \in \mathbb{Z}.
\]

\section*{Notes on Implementation}

The implementation in \texttt{current\_heuristic.py} builds this model as a mixed-integer linear program with two integer variables, using:
\begin{itemize}
\item \texttt{x = LpVariable("buses", lowBound=0, upBound=num\_buses, cat='Integer')}
\item \texttt{y = LpVariable("minibuses", lowBound=0, upBound=num\_minibuses, cat='Integer')}
\end{itemize}

The objective and constraints are added exactly as:
\[
\min \; \texttt{bus\_rental\_cost}\cdot x + \texttt{minibus\_rental\_cost}\cdot y,
\]
subject to
\[
\texttt{bus\_capacity}\cdot x + \texttt{minibus\_capacity}\cdot y \ge \texttt{num\_students},
\]
\[
x+y \le \texttt{num\_drivers}.
\]

The model is solved by the external MILP solver invoked through
\[
\texttt{prob.solve(pulp.GUROBI\_CMD(msg=0))}.
\]
No additional valid inequalities, reformulations, or search logic are present in the code.
